\documentclass{article}

% Language setting
% Replace `english' with e.g. `spanish' to change the document language
\usepackage[english]{babel}

% Set page size and margins
% Replace `letterpaper' with `a4paper' for UK/EU standard size
\usepackage[letterpaper,top=2cm,bottom=2cm,left=3cm,right=3cm,marginparwidth=1.75cm]{geometry}

% Useful packages
\usepackage{amsmath}
\usepackage{amssymb}
\usepackage{graphicx}
\usepackage{inconsolata}
\usepackage{minted}
\usepackage[colorlinks=true, allcolors=blue]{hyperref}

\title{Det ultimate fysikk 1 dokumentet}
\author{Skrevet av André Hansen}

\begin{document}
\maketitle

\tableofcontents

\begin{abstract}
    Dette er et dokument med alle formler og likninger man trenger i fysikk 1. 
    Alt samlet på et sted!
\end{abstract}

\section{Krefter}

\textbf{Kraft $\vec{F}$}

\begin{itemize}
    \item Symbol: $\vec{F}$
    \item Oppgis i Newton($N$)
    \item Kan endre farten til en ting som har en masse
    \item $1N$ tilsvarer å akselerere noe på $1kg$ med $1m/s^2$
    \item Altså newton er masse gange akselerasjon : $kg \cdot m/s^2 = N$
\end{itemize}

\subsection{Newton's lover}

\subsubsection{Newton's 1. lov}

Summen av kreftene er lik null hvis fart er konstant

$$\sum F = 0$$

Denne loven sier egentlig bare at du må ha en kraft for å endre på fart

\subsubsection{Newton's 2. lov}

Krefter på et legeme er massen ganger akselerasjonen

$$\sum F = ma$$

Denne loven deginerer kraft som endring i fart

\subsubsection{Newton's 3. lov}

WIP

\subsubsection{Luftmotsand}

\subsubsection{Friksjonskraft}

\section{Mekanikk}

\subsection{Energi}

\subsubsection{Arbeid}

\begin{itemize}
    \item Et arbeid er bare når noe gir noe annet energi.
    \item Symbol: $\vec{W}$
    \item Oppgis i Newton($J$)
    \item $1J$ tilsvarer å flytte en boks på $1kg$ en meter
    \item Så joules er bare kraft over strekning: $1N \cdot 1m = 1J$
\end{itemize}

\subsubsection{Energi}

Energi er summen av kinetisk energi og potensiell energi

Symbol: $E$

Enhet: $J$

$$E=E_p + E_k$$

\textbf{Potensiell energi: $E_p$}

Potensiell energi er mengen energi tingen kan gjøre til kinetisk energi.
En ball som ligger høyt oppe vil kunne ha potensielle til ha ha høy kinetisk energi fordi den kan rulle ned

$$E_p = mgh$$

\textbf{Kinetisk energi: $E_k$}

Energien på grunn av den beveger seg

$$E_k = \frac{1}{2}mv^2$$

\subsubsection{}

\section{Termofysikk}


\bibliographystyle{alpha}
\bibliography{sample}

\end{document}
